\section*{Capítulo 3: El espaciotiempo, la curvatura y la gravedad}
\begin{description}[style=multiline,leftmargin=2em,labelwidth=1.5em]

% DOCX paragraph 2023
\item[1.]\hypertarget{note:03-01}{}Para quienes estén interesados en una definición matemática
más rigurosa, el espaciotiempo se define como una \emph{variedad} cuyos elementos son sucesos.
A esta variedad se le puede asociar un «mapa» de coordenadas\index{coordenadas} con el cual
se pueden medir las distancias invariantes entre sucesos. Este mapa lo describe un tensor
de rango dos, llamado \emph{tensor métrico}\index{tensor!metrico@métrico} o simplemente
«\emph{métrica}»\index{metrica@métrica|see{tensor, métrico}}. Por tanto, la distancia entre
sucesos se mide mediante el producto escalar de la separación entre sucesos, con la mediación
de la métrica, y es un invariante escalar, es decir, es igual en todos los sistemas de coordenadas.

% DOCX paragraph 2024
\item[2.]\hypertarget{note:03-02}{}Para ser más precisos, la \emph{m} en la ecuación~\eqref{eq:03-02}
indica la llamada \emph{masa en reposo}\index{masa!en reposo}, o sea, la porción de la masa que
no depende del estado de movimiento del objeto. Esta masa en reposo no tiene valor en el caso
de un fotón, cuya masa en reposo es, por tanto, cero. Una expresión más general de la
ecuación~\eqref{eq:03-02}, válida para todo tipo de partículas, es
\emph{E}\textsuperscript{2} \emph{= m}\textsuperscript{2}
\emph{c}\textsuperscript{4} + \emph{p}\textsuperscript{2} \emph{c}\textsuperscript{2},
donde \emph{p} es el \emph{momento lineal} del objeto. Esta expresión también es válida para
un fotón, que siempre tiene un momento no nulo proporcional a su frecuencia.

% DOCX paragraph 2025
\item[3.]\hypertarget{note:03-03}{}Para los lectores interesados, puedo añadir brevemente que
las ecuaciones de Einstein resumidas en la ecuación~\eqref{eq:03-03} son diez ecuaciones
diferenciales\index{ecuaciones!diferenciales} en derivadas parciales de segundo orden y fuertemente
no lineales. Se escriben en forma covariante\index{Einstein!forma covariante} mediante
el uso de dos tensores de rango dos, a saber, el tensor de Einstein\index{Einstein!tensor}
\emph{G\textsubscript{µν}} y el tensor de energía-momento \emph{T\textsubscript{µν}}.

% DOCX paragraph 2026
\item[4.]\hypertarget{note:03-04}{}Nótese que el ejemplo de la sábana y la bola de bolos tiene
un fallo lógico: describe la gravedad mediante el efecto de la bola, pero esta última ya
está sometida a un campo gravitatorio. Más adelante en este capítulo seguiremos explorando
la cuestión, demostrando cómo se puede medir la curvatura sin ni siquiera necesitar una
«sábana». Por el momento basta con señalar que esta analogía es perfectamente válida.
Es importante comprender que la presencia de materia o energía produce curvatura.

% DOCX paragraph 2027
\item[5.]\hypertarget{note:03-05}{}Si el lector o la lectora se anima a hacer un pequeño test,
que preste atención: dados dos puntos A y B sobre la superficie de la Tierra, ¿qué trayectoria
entre A y B minimiza o maximiza la distancia? La respuesta es el círculo máximo que los une:
el círculo que pasa por A, B y el centro de la Tierra. Por ejemplo, si imaginamos que los
dos puntos tienen la misma longitud, entonces la trayectoria seguiría el meridiano que los
une y que pasa por los polos Norte y Sur. Pero, por supuesto, en la mayoría de los casos,
atravesar los polos maximiza la distancia…

\end{description}
